Problem: For what real values of $x$ is $$\sqrt{x + \sqrt{2x-1}} + \sqrt{x - \sqrt{2x-1}} = A$$ given 

\begin{enumerate}
\item $A = \sqrt{2}$
\item $A = 1$
\item $A = 2$
\end{enumerate}

where only non-negative real numbers are admitted for square roots? (Source: 1959 IMO)

Solution: We can actually simplify this equation by squaring both sides. Keep in mind that squaring both sides of an equation is a valid algebraic operation.

\begin{align*}
\sqrt{x + \sqrt{2x-1}} + \sqrt{x - \sqrt{2x-1}} &= A \\
\left(\sqrt{x + \sqrt{2x-1}} + \sqrt{x - \sqrt{2x-1}}\right)\left(\sqrt{x + \sqrt{2x-1}} + \sqrt{x - \sqrt{2x-1}}\right) &= A^2 \\
x + \sqrt{2x-1} + \sqrt{x^2 - (2x - 1)} + \sqrt{x^2 - (2x - 1)} + x - \sqrt{2x-1} = A^2 \\
x + \sqrt{x^2 - (2x - 1)} + \sqrt{x^2 - (2x - 1)} + x = A^2 \\
x + \sqrt{x^2 - 2x + 1} + \sqrt{x^2 - 2x + 1} + x = A^2 \\
x + \sqrt{(x-1)^2} + \sqrt{(x-1)^2} + x = A^2 \\
x + (x-1) + (x-1) + x = A^2 \\
4x - 2 &= A^2 \\
4x &= A^2 + 2 \\
x &= \frac{A^2 + 2}{4}
\end{align*}

Now we can solve for $x$ by plugging in different values of $A$.

When $A = \sqrt{2}$, we have $x = 1$. When $A = 1$, we have $x = \dfrac{3}{4}$. When $A = 2$, we have $x = \dfrac{3}{2}$.

Thus $x = \left\{1, \dfrac{3}{4}, \dfrac{3}{2} \right\}$.
